\import{relation/properties.tex}
\import{relation/equivalence.tex}

\section{Closure operations on relations}

\begin{definition}\label{reflexive_closure}
    $\reflexiveClosure{X}{R} = R\union\identity{X}$.
\end{definition}

\begin{classification}\label{reflexive_closure_binary_relation_on}
    Suppose $R$ is a binary relation on $X$.
    Then $\reflexiveClosure{X}{R}$ is a binary relation on $X$.
\end{classification}
\begin{proof}
    $R\subseteq X\times X$ by definition.
    $\identity{X}\subseteq X\times X$
        by \obv{identity_binary_relation_on}.
    Hence $\reflexiveClosure{X}{R}\subseteq X\times X$ by
        \auto{reflexive_closure,subseteq,union_iff}.
    Follows by definition.
\end{proof}

Adjoining the identity relation on $X$ makes a relation reflexive on $X$.

\begin{proposition}\label{reflexive_closure_is_reflexive}
    $\reflexiveClosure{X}{R}$ is reflexive on $X$.
\end{proposition}
\begin{proof}
    Follows by \auto{reflexive_closure,reflexive_on,union_iff,id_iff}.
\end{proof}

\begin{proposition}\label{converse_reflexive_closure}
    $\converse{\reflexiveClosure{X}{R}}
        =\reflexiveClosure{X}{\converse{R}}$.
\end{proposition}
\begin{proof}
    Follows by
        \auto{reflexive_closure,converse_union,converse_identity}.
\end{proof}

\begin{definition}\label{reflexive_reduction}
    $\reflexiveReduction{X}{R} = R\setminus\identity{X}$.
\end{definition}

Symmetric closure adds the converse relation. Transitive and
reflexive-transitive closures are developed in the following chapter; an
equivalence-closure construction remains deferred.

\begin{definition}\label{symmetric_closure}
    $\symmetricClosure{R} = R\union\converse{R}$.
\end{definition}
